% Generated draft: requires mathematical and source review.
\documentclass[11pt]{article}
\usepackage[margin=25mm]{geometry}
\usepackage{fontspec}
\setmainfont{Latin Modern Roman}
\usepackage{amsmath,amssymb,mathtools,unicode-math}
\setmathfont{Latin Modern Math}
\usepackage{enumitem,needspace,xurl,hyperref}
\setlength{\parindent}{0pt}
\setlength{\parskip}{5pt}
\setlength{\emergencystretch}{3em}
\newcommand{\N}{\mathbb{N}}
\newcommand{\Z}{\mathbb{Z}}
\newcommand{\Q}{\mathbb{Q}}
\newcommand{\R}{\mathbb{R}}
\newcommand{\C}{\mathbb{C}}
\newcommand{\F}{\mathbb{F}}
\newcommand{\A}{\mathbb{A}}
\newcommand{\PP}{\mathbb{P}}
\newcommand{\E}{\mathbb{E}}
\newcommand{\eps}{\varepsilon}
\newcommand{\dd}{\,\mathrm{d}}
\begin{document}
\section*{UnsolvedMath ID: 800287}
\textbf{Source problem code:} catalogue entry
\subsection{Definitions and mathematical statement}
Let \(I\subseteq\R\) denote the nondegenerate interval intended in the source, with its usual order and topology. Write \(\operatorname{Homeo}_{+}(I)\) for the group, under composition, of strictly increasing homeomorphisms \(I\to I\). An orientation-preserving action of a group \(\Gamma\) on \(I\) is a homomorphism \(\rho:\Gamma\to\operatorname{Homeo}_{+}(I)\). It is nontrivial if
\[
\exists\gamma\in\Gamma\;\exists x\in I\quad \rho(\gamma)(x)\ne x.
\]
In particular, faithfulness is not required. A lattice in a Lie group \(G\) is a discrete subgroup \(\Gamma\leq G\) such that \(G/\Gamma\) admits a nonzero finite \(G\)-invariant measure.

The question attributed to Navas (2017) is: does there exist a higher-rank simple Lie group \(G\), a lattice \(\Gamma\leq G\), and a homomorphism
\[
\rho:\Gamma\longrightarrow\operatorname{Homeo}_{+}(I)
\]
whose image is not the trivial group?

The supplied record does not specify the interval or its endpoint convention. For a closed interval, every orientation-preserving homeomorphism fixes both endpoints; nontriviality therefore requires some interior point to move. The record also does not explicitly specify the rank convention or the global hypotheses on \(G\), including connectedness and restrictions on its center. Under the usual real-rank interpretation, the higher-rank condition is \(\operatorname{rank}_{\R}(G)\geq2\). These conventions require confirmation before the question has a fully fixed domain.
\subsection{Short English statement}
Navas asks whether some lattice in a higher-rank simple Lie group has an orientation-preserving action on an interval in which at least one group element moves a point. The supplied registry labels the problem open; this label is not independently verified. The intended interval and precise Lie-group conventions remain to be confirmed.
\subsection{Sources}
\begin{itemize}
\item Local UnsolvedMath registry, permanent ID \(800287\), source problem code \texttt{catalogue entry}, title: \emph{Group Actions on 1-Manifolds --- Question 1}. Supplied attribution: Navas, 2017. Supplied status: open.
\item Navas, \emph{Group Actions on 1-Manifolds: A List of Very Concrete Open Questions} (2017), Question 1 (\#1). Supplied source URL: \url{https://arxiv.org/abs/1712.06462}. The provenance records access on 2026-10-16 and extraction by \texttt{arxiv\_latex}; these are supplied provenance claims, not an independent source visit or verification.
\item Supplied upstream revision: \texttt{d371ee917ac9647fbb06cd6fc529deec6f3dda82}. The registry record was created and updated on 2026-10-16.
\end{itemize}
\end{document}
